% ═══════════════════════════════════════════════════════════════
% SPANISCH LEHRERHINWEISE - Sequenz 6a: Los lugares
% ═══════════════════════════════════════════════════════════════
% Fachfarbe: #c41e3a (Karminrot)
% ═══════════════════════════════════════════════════════════════

\documentclass[11pt, a4paper]{article}

% SW-MODUS TOGGLE
\newif\ifswmodus
\swmodusfalse

% PAKETE
\usepackage[utf8]{inputenc}
\usepackage[T1]{fontenc}
\usepackage[ngerman]{babel}
\usepackage{helvet}
\renewcommand{\familydefault}{\sfdefault}

\usepackage[margin=2cm]{geometry}
\usepackage{enumitem}
\usepackage{tabularx}
\usepackage{booktabs}
\usepackage{longtable}

\usepackage{xcolor}
\usepackage{amssymb}
\usepackage{tcolorbox}
\tcbuselibrary{skins}

\usepackage{fancyhdr}
\usepackage{lastpage}
\usepackage{needspace}

% FARBEN
\definecolor{spanisch-color}{HTML}{c41e3a}
\definecolor{grau-color}{HTML}{888888}
\definecolor{hellgrau-color}{HTML}{f5f5f5}
\definecolor{basis-color}{HTML}{2a7a4b}
\definecolor{standard-color}{HTML}{2c5aa0}
\definecolor{erweiterung-color}{HTML}{7b1fa2}
\definecolor{loesung-color}{HTML}{c62828}

\definecolor{sw-dark}{HTML}{000000}
\definecolor{sw-gray}{HTML}{666666}
\definecolor{sw-white}{HTML}{ffffff}

\ifswmodus
    \colorlet{spanisch}{sw-dark}
    \colorlet{grau}{sw-gray}
    \colorlet{hellgrau}{sw-white}
    \colorlet{basis}{sw-dark}
    \colorlet{standard}{sw-dark}
    \colorlet{erweiterung}{sw-dark}
    \colorlet{loesung}{sw-dark}
\else
    \colorlet{spanisch}{spanisch-color}
    \colorlet{grau}{grau-color}
    \colorlet{hellgrau}{hellgrau-color}
    \colorlet{basis}{basis-color}
    \colorlet{standard}{standard-color}
    \colorlet{erweiterung}{erweiterung-color}
    \colorlet{loesung}{loesung-color}
\fi

\newcommand{\fachfarbe}{spanisch}

% VARIABLEN
\newcommand{\thema}{Los lugares -- Orte in der Stadt}
\newcommand{\sequenz}{06}
\newcommand{\block}{a}
\newcommand{\fach}{Spanisch}
\newcommand{\klasse}{7}
\newcommand{\dauer}{90 Min. (Doppelstunde)}
\newcommand{\lerngruppengroesse}{25}
\newcommand{\datum}{\today}

% KOPF- UND FUSSZEILE
\pagestyle{fancy}
\fancyhf{}
\renewcommand{\headrulewidth}{0.4pt}
\renewcommand{\footrulewidth}{0.4pt}

\lhead{\fach{} -- Klasse \klasse}
\chead{\textbf{LH-\sequenz\block{} (Lehrerhinweise)}}
\rhead{\datum}

\lfoot{}
\cfoot{}
\rfoot{Seite \thepage\ von \pageref{LastPage}}

% BEFEHLE
\newcommand{\B}{\textcolor{basis}{\textbf{[B]}}}
\newcommand{\Std}{\textcolor{standard}{\textbf{[S]}}}
\newcommand{\E}{\textcolor{erweiterung}{\textbf{[E]}}}
\newcommand{\loesung}[1]{\textcolor{loesung}{\textbf{#1}}}
\newcommand{\es}[1]{\textcolor{\fachfarbe}{\textbf{#1}}}

\newtcolorbox{kurzplan}{
    colback=hellgrau,
    colframe=\fachfarbe,
    colbacktitle=white,
    coltitle=\fachfarbe,
    boxrule=1pt,
    arc=0pt,
    fonttitle=\bfseries,
    attach boxed title to top left={yshift=-2mm, xshift=5mm},
    boxed title style={boxrule=0pt, colback=white},
    title=Kurzplan
}

\newtcolorbox{differenzierung}{
    colback=white,
    colframe=grau,
    colbacktitle=white,
    coltitle=grau,
    boxrule=0.5pt,
    arc=0pt,
    fonttitle=\bfseries,
    attach boxed title to top left={yshift=-2mm, xshift=5mm},
    boxed title style={boxrule=0pt, colback=white},
    title=Differenzierung
}

\newtcolorbox{fehlerbox}{
    colback=white,
    colframe=loesung,
    colbacktitle=white,
    coltitle=loesung,
    boxrule=0.5pt,
    arc=0pt,
    fonttitle=\bfseries,
    attach boxed title to top left={yshift=-2mm, xshift=5mm},
    boxed title style={boxrule=0pt, colback=white},
    title=Haufige Fehler
}

\newtcolorbox{materialbox}{
    colback=white,
    colframe=\fachfarbe,
    colbacktitle=white,
    coltitle=\fachfarbe,
    boxrule=0.5pt,
    arc=0pt,
    fonttitle=\bfseries,
    attach boxed title to top left={yshift=-2mm, xshift=5mm},
    boxed title style={boxrule=0pt, colback=white},
    title=Material-Checkliste
}

% ═══════════════════════════════════════════════════════════════
% DOKUMENT
% ═══════════════════════════════════════════════════════════════

\begin{document}

% KOPFBEREICH
\begin{center}
    {\Large\bfseries\textcolor{\fachfarbe}{Lehrerhinweise: \thema}}\\[0.3em]
    {\normalsize LH-\sequenz\block{} | \dauer}
\end{center}

\vspace{10pt}

% ═══════════════════════════════════════════════════════════════
% 1. KURZPLAN
% ═══════════════════════════════════════════════════════════════

\section*{1. Stundenverlauf (Kurzplan)}

\textbf{1. Stunde (45 Min.) -- Orte kennenlernen}

\begin{tabularx}{\textwidth}{|c|l|X|l|}
\hline
\textbf{Zeit} & \textbf{Phase} & \textbf{Inhalt} & \textbf{Material} \\
\hline
0--3 & Einstieg & Begrußung, Stadtbild zeigen & Bild \\
\hline
3--8 & Impuls & Was gibt es in einer Stadt? Vorwissen sammeln & Tafel \\
\hline
8--18 & Input & 10 Orte mit Bildkarten prasentieren, Chor & Bildkarten \\
\hline
18--28 & Sortieren & PA: Nach Artikel sortieren (el/la) & AB \\
\hline
28--35 & Austausch & Ergebnisse zusammentragen, Genus-Regel & Tafel \\
\hline
35--42 & Merkkasten & SuS ubertragen Vokabeln in AB & AB \\
\hline
42--45 & Ausblick & Verabschiedung, Ausblick nachste Stunde & -- \\
\hline
\end{tabularx}

\vspace{1em}

\textbf{2. Stunde (45 Min.) -- Orte anwenden}

\begin{tabularx}{\textwidth}{|c|l|X|l|}
\hline
\textbf{Zeit} & \textbf{Phase} & \textbf{Inhalt} & \textbf{Material} \\
\hline
0--3 & Begrußung & Wiederholung: Blitzabfrage Orte & -- \\
\hline
3--10 & Quiz & Bildkarten-Quiz: Bild zeigen, Ort nennen & Bildkarten \\
\hline
10--20 & Funktionen & PA: Aktivitaten zu Orten zuordnen & AB \\
\hline
20--32 & Stadtplan & PA/GA: Orte im Stadtplan lokalisieren & Stadtplan \\
\hline
32--42 & AB-Arbeit & EA: AB-06a Aufgaben 2--4 & AB \\
\hline
42--47 & Besprechung & Losungen vergleichen, Fehler besprechen & -- \\
\hline
47--50 & Abschluss & HA, Verabschiedung & -- \\
\hline
\end{tabularx}

% ═══════════════════════════════════════════════════════════════
% 2. DIFFERENZIERUNG
% ═══════════════════════════════════════════════════════════════

\section*{2. Differenzierung}

\begin{differenzierung}
\textbf{Legende:} \B{} Basis (BR) | \Std{} Standard (MR) | \E{} Erweiterung (GY)

\vspace{0.5em}

\textbf{WICHTIG:} Diese Marker sind NUR fur die Lehrkraft. Auf Schulermaterial erscheinen KEINE Niveaumarkierungen!
\end{differenzierung}

\vspace{1em}

\begin{tabularx}{\textwidth}{|l|X|X|X|}
\hline
& \B{} \textbf{Basis} & \Std{} \textbf{Standard} & \E{} \textbf{Erweiterung} \\
\hline
\textbf{Vokabeln} & 6 Kernorte (parque, cine, hospital, supermercado, farmacia, estación) & Alle 10 Orte & + teatro, ayuntamiento, iglesia \\
\hline
\textbf{Aufgabe 1} & Mit Bildunterstutzung & Selbststandig & Zusatz: Satze bilden \\
\hline
\textbf{Aufgabe 2} & Artikel vorgeben, nur ankreuzen & Selbststandig einsetzen & + Begrundung fur Genus \\
\hline
\textbf{Aufgabe 3} & Wortliste mit Bildern & Nur Wortliste & Ohne Wortliste \\
\hline
\textbf{Aufgabe 4} & 2 Satze mit Muster & 3 Satze & 5 Satze, freie Formulierung \\
\hline
\end{tabularx}

\vspace{1em}

\textbf{Konkrete Hinweise:}
\begin{itemize}
    \item \B{} SuS mit LRS: Zeitzugabe (+10 Min.), vergroßerte AB-Version, Bildkarten als Hilfe
    \item \B{} DaZ-SuS: Wortschatz deutsch-spanisch vorab besprechen, ggf. auch auf Englisch
    \item \E{} Leistungsstarke: Zusatzorte lernen, Expertenrolle bei Partnerarbeit, Stadtbeschreibung schreiben
\end{itemize}

% ═══════════════════════════════════════════════════════════════
% 3. HAUFIGE FEHLER
% ═══════════════════════════════════════════════════════════════

\section*{3. Haufige Fehler}

\begin{fehlerbox}
\begin{tabularx}{\textwidth}{|l|X|X|}
\hline
\textbf{Fehler} & \textbf{Beispiel} & \textbf{Reaktion} \\
\hline
Falscher Artikel & *\es{la cine} statt \es{el cine} & Regel: -o meist maskulin, -a meist feminin \\
\hline
Aussprache biblioteca & *[bi-bli-o-te-ka] mit 5 Silben & Langsam: bi-blio-te-ca (4 Silben) \\
\hline
Orthographie farmacia & *pharmacia, *farmazia & Vergleich EN: pharmacy vs. ES: farmacia \\
\hline
Akzent estación & *estacion ohne Akzent & Betonung auf letzter Silbe = Akzent \\
\hline
Genus centro comercial & *la centro comercial & Centro = maskulin, bestimmt den Artikel \\
\hline
\end{tabularx}
\end{fehlerbox}

% ═══════════════════════════════════════════════════════════════
% 4. MATERIAL-CHECKLISTE
% ═══════════════════════════════════════════════════════════════

\section*{4. Material-Checkliste}

\begin{materialbox}
\begin{itemize}[label=$\square$]
    \item AB-\sequenz\block.pdf (\lerngruppengroesse{} Kopien)
    \item ML-\sequenz\block.pdf (1$\times$ fur Lehrkraft)
    \item Lehrwerk: Que pasa 1, Unidad 6
    \item Bildkarten: 10 Orte (A4, laminiert empfohlen)
    \item Stadtplan ``Villa Ejemplo'' (Fantasiestadt, \lerngruppengroesse{} Kopien)
    \item Stadtbild Madrid/Barcelona (Beamer oder A3)
    \item Optional: LP-\sequenz\block.html (Tablets/PCs)
\end{itemize}
\end{materialbox}

% ═══════════════════════════════════════════════════════════════
% 5. LOSUNGEN
% ═══════════════════════════════════════════════════════════════

\section*{5. Losungen AB-\sequenz\block}

\textbf{Merkkasten: Orte in der Stadt}

\begin{tabular}{ll}
\es{el parque} & $\rightarrow$ \loesung{der Park} \\
\es{la plaza} & $\rightarrow$ \loesung{der Platz} \\
\es{el cine} & $\rightarrow$ \loesung{das Kino} \\
\es{la biblioteca} & $\rightarrow$ \loesung{die Bibliothek} \\
\es{el museo} & $\rightarrow$ \loesung{das Museum} \\
\es{la farmacia} & $\rightarrow$ \loesung{die Apotheke} \\
\es{el supermercado} & $\rightarrow$ \loesung{der Supermarkt} \\
\es{la estación} & $\rightarrow$ \loesung{der Bahnhof} \\
\es{el hospital} & $\rightarrow$ \loesung{das Krankenhaus} \\
\es{el centro comercial} & $\rightarrow$ \loesung{das Einkaufszentrum} \\
\end{tabular}

\vspace{1em}

\textbf{Aufgabe 1: Zuordnung} (5 Punkte)

\begin{tabular}{ll}
\es{el parque} & $\rightarrow$ \loesung{B) der Park} \\
\es{el cine} & $\rightarrow$ \loesung{D) das Kino} \\
\es{la farmacia} & $\rightarrow$ \loesung{A) die Apotheke} \\
\es{la estación} & $\rightarrow$ \loesung{C) der Bahnhof} \\
\es{el hospital} & $\rightarrow$ \loesung{E) das Krankenhaus} \\
\end{tabular}

\vspace{1em}

\textbf{Aufgabe 2: Artikel einsetzen} (5 Punkte)

\begin{tabular}{ll}
a) & \loesung{el} museo \\
b) & \loesung{la} plaza \\
c) & \loesung{el} supermercado \\
d) & \loesung{la} biblioteca \\
e) & \loesung{el} centro comercial \\
\end{tabular}

\vspace{1em}

\textbf{Aufgabe 3: Luckentext} (5 Punkte)

\begin{tabular}{ll}
a) & \loesung{la biblioteca} \\
b) & \loesung{el supermercado} \\
c) & \loesung{el cine} \\
d) & \loesung{el hospital} \\
e) & \loesung{el parque} \\
\end{tabular}

\vspace{1em}

\textbf{Aufgabe 4: Satze schreiben} (5 Punkte)

Mogliche Losungen:
\begin{itemize}
    \item \loesung{En mi ciudad hay un cine.}
    \item \loesung{La biblioteca está cerca del parque.}
    \item \loesung{El hospital es grande.}
    \item \loesung{Me gusta el parque.}
    \item \loesung{El supermercado está en el centro.}
\end{itemize}

\textit{Bewertung: Ort korrekt (1P), Artikel korrekt (0,5P), Satzstruktur korrekt (0,5P) = max. 2P pro Satz, bei 3 Satzen = max. 5P (leichte Aufwertung fur alle 3)}

% ═══════════════════════════════════════════════════════════════
% 6. HAUSAUFGABE
% ═══════════════════════════════════════════════════════════════

\section*{6. Hausaufgabe}

\begin{itemize}
    \item Vokabeln lernen: 10 Orte mit Artikel (Merkkasten)
    \item Optional: Lernpfad LP-\sequenz\block.html (digital)
\end{itemize}

% ═══════════════════════════════════════════════════════════════
% 7. REFLEXIONSNOTIZEN
% ═══════════════════════════════════════════════════════════════

\section*{7. Reflexionsnotizen (nach Durchfuhrung)}

\begin{tabularx}{\textwidth}{|l|X|}
\hline
\textbf{Was lief gut?} & \\[2em]
\hline
\textbf{Was lief nicht?} & \\[2em]
\hline
\textbf{Zeitmanagement} & \\[2em]
\hline
\textbf{Fur nachstes Mal} & \\[2em]
\hline
\end{tabularx}

\end{document}
